\documentclass[twocolumn]{aastex631}

\usepackage{amsmath}

\newcommand{\tic}{TIC~402710280}
\newcommand{\planet}{TIC~402710280.01}
\newcommand{\tess}{\textit{TESS}}
\newcommand{\gaia}{\textit{Gaia}}
\newcommand{\rearth}{R_\oplus}
\newcommand{\rsun}{R_\odot}
\newcommand{\msun}{M_\odot}

\shorttitle{A 36.34-day sub-Neptune candidate around TIC 402710280}
\shortauthors{Vora}

\begin{document}

\title{TIC 402710280.01: A Sub-Neptune Candidate on a 36.34-Day Orbit around a Nearby K Dwarf, Identified in TESS Full-Frame Images}

\author[0000-0000-0000-0000]{Rishit Vora}
\affiliation{Independent Researcher}
\email{EMAIL}

\begin{abstract}
We report a transiting planet candidate around \tic\ (TYC~9041-03677-1), a $T = 9.24$ K dwarf at $66.8$~pc that the Transiting Exoplanet Survey Satellite (\tess) has observed only in full-frame images. Four transits, recorded in Sectors~39, 93, 100 and~102 and reduced by three independent pipelines, share a depth of $2097 \pm 46$~ppm and a duration of $5.27 \pm 0.05$~h. Two of the transits fall on consecutive orbits, and the predicted transits of every shorter period are absent from the data at more than $18\sigma$; the orbital period is $P = 36.33990 \pm 0.00053$~d. A transit model with each event centred on its measured mid-time yields $R_p/R_\star = 0.0421^{+0.0006}_{-0.0005}$, impact parameter $b = 0.22^{+0.19}_{-0.15}$ and a stellar density of $1.57^{+0.11}_{-0.29}\,\rho_\odot$, in agreement with the catalogue value of $1.61 \pm 0.38\,\rho_\odot$. With $R_\star = 0.817 \pm 0.047\,\rsun$ the candidate has a radius of $3.75 \pm 0.22\,\rearth$, receives $10.2 \pm 1.7$ times Earth's insolation and has an equilibrium temperature of $498 \pm 21$~K. Difference-image centroids locate the signal within $1.5'' \pm 1.3''$ of the target; the only catalogued star that could produce the signal, a background source $3.5''$ away, is disfavoured by the transit shape at $\Delta\chi^2 \geq 142$. No secondary eclipse is detected ($-54 \pm 39$~ppm), and \texttt{TRICERATOPS} returns a false-positive probability of $0.003 \pm 0.004$. The mid-times scatter about a linear ephemeris by $21.5$~min (rms), which we report without interpretation. We publish predicted transit times through 2028, including a transit expected in \tess\ Sector~114 on 2027 February~19.
\end{abstract}

\keywords{Exoplanets (498) --- Transit photometry (1709) --- Mini Neptunes (1063) --- K dwarf stars (876)}

\section{Introduction} \label{sec:intro}

Transiting planets with periods longer than a month around bright, nearby stars remain scarce, because a single \tess\ sector of about 27~d \citep{Ricker2015} records at most one of their transits. Such planets are found by combining sectors observed years apart, a task that becomes tractable as the mission's full-frame images accumulate. Stars observed only in full-frame images are additionally excluded from the Science Processing Operations Center (SPOC) search of two-minute data \citep{Jenkins2016}, so their long-period transiters depend on full-frame light curves such as those of the Quick-Look Pipeline \citep[QLP;][]{Huang2020a, Huang2020b} and TESS-SPOC \citep{Caldwell2020}, and on photometry of the most recent images, which these products do not yet cover.

Here we report \planet, a sub-Neptune candidate with a 36.34-day period around a K dwarf at 67~pc. Its first two transits were identified in public full-frame light curves; the period was determined from our own photometry of the 2026 full-frame images. As of 2026 October~14, the star is neither a \tess\ Object of Interest \citep[TOI;][]{Guerrero2021} nor a community TOI, and no planet parameters, follow-up observations or literature references are associated with it in the Exoplanet Follow-up Observing Program (ExoFOP) archive, the NASA Exoplanet Archive \citep{Akeson2013}, or the NASA Astrophysics Data System.

\section{Host Star} \label{sec:star}

Table~\ref{tab:star} lists the properties of \tic. We adopt the stellar parameters of the \tess\ Input Catalog v8.2 \citep[TIC;][]{Stassun2019}. \gaia\ DR3 \citep{Gaia2016, GaiaDR3} reports a parallax of $14.964 \pm 0.014$~mas, a renormalised unit weight error \citep[RUWE;][]{Lindegren2021} of $1.00$, no multi-peaked image windows (\texttt{ipd\_frac\_multi\_peak} $= 0$), and no non-single-star solution. Together these indicate a single, well-behaved astrometric source, although they do not exclude a companion at small separation and low mass ratio.

\begin{deluxetable}{lcc}
\tablecaption{Properties of \tic \label{tab:star}}
\tablehead{\colhead{Parameter} & \colhead{Value} & \colhead{Source}}
\startdata
R.A. (J2000) & 16:34:05.67 & TIC \\
Decl. (J2000) & $-$63:32:58.8 & TIC \\
\gaia\ DR3 ID & 5828441289812582400 & \gaia \\
$\mu_\alpha \cos\delta$ (mas yr$^{-1}$) & $59.909$ & \gaia \\
$\mu_\delta$ (mas yr$^{-1}$) & $13.685$ & \gaia \\
Parallax (mas) & $14.964 \pm 0.014$ & \gaia \\
Distance (pc) & $66.83 \pm 0.06$ & \gaia \\
RUWE & $1.00$ & \gaia \\
$T$ (mag) & $9.243 \pm 0.006$ & TIC \\
$G$ (mag) & $9.781$ & \gaia \\
$V$ (mag) & $9.982 \pm 0.004$ & TIC \\
$K_s$ (mag) & $8.026$ & TIC \\
$T_{\rm eff}$ (K) & $5182 \pm 121$ & TIC \\
$\log g$ (cgs) & $4.56 \pm 0.08$ & TIC \\
$R_\star$ ($\rsun$) & $0.817 \pm 0.047$ & TIC \\
$M_\star$ ($\msun$) & $0.88 \pm 0.11$ & TIC \\
$\rho_\star$ ($\rho_\odot$) & $1.61 \pm 0.38$ & TIC \\
\enddata
\end{deluxetable}

\section{Observations and Photometry} \label{sec:data}

\tess\ observed \tic\ in Sectors~12 (2019), 39 (2021), 66 (2023), 93 (2025) and 100--103 (2026), always in full-frame images, at cadences of 1800~s (prime mission), 600~s (first extended mission) and 200~s thereafter. We use the TESS-SPOC light curves \citep{Caldwell2020, TESSSPOC} for Sectors~12, 39 and~66 (pre-search data conditioning flux) and the QLP light curve \citep{Huang2020a, QLPdata} for Sector~93 (systematics-corrected flux). Neither pipeline had processed Sectors~100--103 when this work was done, so we extracted their light curves from the calibrated full-frame images: $21 \times 21$-pixel cutouts \citep{Brasseur2019}, a per-cadence background equal to the median of the faintest 40\% of pixels away from the target, a circular aperture whose radius (between 1 and 4 pixels) minimises the one-hour scatter of the detrended light curve, and rejection of cadences flagged by SPOC or affected by scattered-light excursions of the background.

For transit analysis, each event is extracted in a window of $\pm 0.8$~d, cadences with non-zero quality flags are removed, and the flux is normalised by a quadratic function of time fitted to the out-of-transit data. For timing, the baseline is instead refitted simultaneously with the transit model (\S\ref{sec:times}).

\section{Detection and Period} \label{sec:period}

A search of public full-frame light curves for individual transit-like events (box filters at ten trial durations, with thresholds set by the measured scatter of each filtered series) identified two dips of matching depth and duration in Sectors~39 and 93, separated by 1453.58~d. Our photometry of Sectors~99--108 then revealed two further transits, in Sectors~100 and 102, separated by a single orbit of 36.36~d (Figure~\ref{fig:transits}). Because these two transits are consecutive, the period cannot be a multiple of 36.36~d; the remaining possibilities are its integer fractions. We tested $P/2$ through $P/5$ by measuring the flux at every transit time they predict within the data, excluding the four observed events: the shallowest exclusion among the 41 covered epochs is $18.6\sigma$ (for $P/4$), and every alias is rejected. Sector~101 contains no transit at the adopted period, Sector~103's predicted transit falls in a data gap, and Sector~12 contains no predicted transit. The four transits thus fix the period uniquely at 36.34~d.

The transit predicted in Sector~66 (BTJD~3111.40, where BTJD $=$ BJD$_{\rm TDB} - 2457000$) coincides with an interval of strong scattered light (Figure~\ref{fig:s66}). The background falls from $1.5 \times 10^4$ to $2.9 \times 10^3$~e$^-$~s$^{-1}$ across the window, quality-flagged cadences end only 3.2~h before the predicted mid-time (0.6~h before ingress), and QLP flags every cadence within 12~h of it. A transit model fitted at the predicted time returns a depth of $1620 \pm 200$~ppm, consistent with the other events, but the result depends on the treatment of the one-sided baseline. We therefore exclude Sector~66 from all fits; it neither supports nor contradicts the candidate.

\begin{figure*}
\epsscale{1.15}
\plotone{figures/fig1.pdf}
\caption{Transits of \planet. Grey points are individual cadences; blue points are 20-minute bins with their standard errors. Each window is centred on its measured mid-time (Table~\ref{tab:times}), and the red curve is the median transit model of \S\ref{sec:fit}. The right-hand panel shows the Sector~66 window at the time predicted by the linear ephemeris; data with zero quality flags begin at the dotted line, and the window is not used in any fit (\S\ref{sec:period}).\label{fig:transits}}
\end{figure*}

\begin{figure}
\plotone{figures/fig2.pdf}
\caption{TESS-SPOC data of Sector~66 around the predicted transit (shaded). From top: simple-aperture flux, with quality-flagged cadences in grey; pre-search data conditioning flux, quality 0 only; and the background. The window lies on a scattered-light ramp and is not used in the analysis.\label{fig:s66}}
\end{figure}

\section{Transit Analysis} \label{sec:analysis}

\subsection{Individual Transit Times and Depths} \label{sec:times}

Each transit is timed by fixing the shape to the median model of \S\ref{sec:fit} and scanning the mid-time on a 0.5-minute grid over $\pm 2$~h, refitting a quadratic baseline at each step. The $1\sigma$ uncertainty is the half-width of the $\Delta\chi^2 = 1$ interval, multiplied by the red-noise factor $\beta$ \citep{Pont2006, Winn2008}, the median ratio of the binned residual scatter (10--60-minute bins) to its white-noise expectation; $\beta$ ranges from 1.00 to 1.53. The depth of each event is measured with the shape fixed and a free amplitude. Table~\ref{tab:times} lists the results.

The four depths agree ($\chi^2 = 2.3$ for 3 degrees of freedom) despite coming from three independent reductions with different apertures and cadences, which argues against an origin in any one pipeline. Their weighted mean is $2097 \pm 46$~ppm.

\begin{deluxetable*}{lcccccc}
\tablecaption{Individual transits of \planet \label{tab:times}}
\tablehead{\colhead{Sector} & \colhead{Light curve} & \colhead{Cadence (s)} & \colhead{Epoch} & \colhead{$T_c$ (BJD$_{\rm TDB}$)} & \colhead{$O - C$ (min)} & \colhead{Depth (ppm)}}
\startdata
39 & TESS-SPOC & 600 & 0 & $2459384.6094 \pm 0.0013$ & $+4.8$ & $2203 \pm 85$ \\
93 & QLP & 200 & 40 & $2460838.1847 \pm 0.0012$ & $-24.8$ & $2050 \pm 82$ \\
100 & This work & 200 & 47 & $2461092.5851 \pm 0.0024$ & $+5.7$ & $2075 \pm 150$ \\
102 & This work & 200 & 48 & $2461128.9449 \pm 0.0016$ & $+34.3$ & $2045 \pm 86$ \\
\enddata
\tablecomments{$O-C$ is relative to the linear ephemeris of Table~\ref{tab:params}. Depths are the model depth at mid-transit for a free amplitude with the shape fixed.}
\end{deluxetable*}

\subsection{Transit Shape} \label{sec:fit}

We model the four transits jointly with \texttt{batman} \citep{Kreidberg2015}, integrating over each light curve's exposure time and centring each window on its measured mid-time, so that timing scatter (\S\ref{sec:ttv}) does not smear the stacked profile. The free parameters are $R_p/R_\star$, $b$ and $\log_{10}\rho_\star$, with uniform priors; $a/R_\star$ follows from $\rho_\star$ and $P$ \citep{Seager2003} for a circular orbit. Quadratic limb-darkening coefficients are fixed at $(u_1, u_2) = (0.39, 0.21)$, interpolated in $T_{\rm eff}$ from the \tess-band tables of \citet{Claret2017} at $\log g = 4.5$. Each window's flux uncertainty is its out-of-transit point-to-point scatter multiplied by its $\beta$. We sample the posterior with \texttt{emcee} \citep{Goodman2010, ForemanMackey2013}, using 40 walkers for 30\,000 steps and discarding the first third; the chain spans more than 100 integrated autocorrelation times for every parameter. Table~\ref{tab:params} gives the medians and 68\% intervals.

The fitted stellar density, $1.57^{+0.11}_{-0.29}\,\rho_\odot$, agrees with the TIC value of $1.61 \pm 0.38\,\rho_\odot$ within $0.1\sigma$; the transit is therefore consistent with a planet on a circular orbit around \tic\ itself. The best-fitting model has $\chi^2 = 1866$ for 2227 data points.

\subsection{Ephemeris and Timing Residuals} \label{sec:ttv}

A weighted linear fit to the four mid-times gives $P = 36.33990$~d and a reference epoch $T_0 = 2460838.2019$~BJD$_{\rm TDB}$, placed at the Sector~93 transit where the covariance between $T_0$ and $P$ is smallest. The residuals (Table~\ref{tab:times}, Figure~\ref{fig:oc}) are $+4.8$, $-24.8$, $+5.7$ and $+34.3$~min, with an rms of 21.5~min, much larger than the 1.7--3.5-min formal uncertainties ($\chi^2 = 461$ for 2 degrees of freedom). We scale the ephemeris uncertainties by $\sqrt{\chi^2/\nu}$, which yields $\sigma_P = 0.00053$~d and $\sigma_{T_0} = 0.0122$~d.

The scatter has two possible origins: timing systematics that the formal uncertainties do not capture, or transit-timing variations induced by an additional body. The largest residuals belong to Sector~93, a public QLP light curve, and Sector~102, whose transit occurs 1.4~d after the start of an orbit, where a residual ramp is visible; no single reduction is therefore responsible. Four epochs cannot distinguish the two explanations, and we draw no conclusion about transit-timing variations. The predictions in \S\ref{sec:predictions} include the observed scatter.

\begin{figure}
\plotone{figures/fig3.pdf}
\caption{Transit-timing residuals relative to the linear ephemeris. Symbols denote the light-curve source: TESS-SPOC (square), QLP (diamond) and our photometry (open circles). The shaded band is the $1\sigma$ uncertainty of the ephemeris; the dotted curves add the 21.5-min rms scatter of the residuals in quadrature. Dashed lines mark epochs 57 (2027 February~19, \tess\ Sector~114) and 59 (2027 May~3).\label{fig:oc}}
\end{figure}

\begin{deluxetable}{lc}
\tablecaption{Parameters of \planet \label{tab:params}}
\tablehead{\colhead{Parameter} & \colhead{Value}}
\startdata
\multicolumn{2}{c}{\textit{Ephemeris}} \\
$P$ (d) & $36.33990 \pm 0.00053$ \\
$T_0$ (BJD$_{\rm TDB}$) & $2460838.2019 \pm 0.0122$ \\
Timing residual rms (min) & 21.5 \\
\multicolumn{2}{c}{\textit{Fitted}} \\
$R_p/R_\star$ & $0.04208^{+0.00064}_{-0.00051}$ \\
$b$ & $0.22^{+0.19}_{-0.15}$ \\
$\rho_\star$ ($\rho_\odot$) & $1.57^{+0.11}_{-0.29}$ \\
\multicolumn{2}{c}{\textit{Derived}} \\
Depth at mid-transit (ppm) & $2097 \pm 46$ \\
$T_{14}$ (h) & $5.27^{+0.05}_{-0.04}$ \\
$a/R_\star$ & $53.7^{+1.2}_{-3.5}$ \\
$R_p$ ($\rearth$) & $3.75 \pm 0.22$ \\
$a$ (AU) & $0.206 \pm 0.008$ \\
$S$ ($S_\oplus$) & $10.2 \pm 1.7$ \\
$T_{\rm eq}$ (K) & $498 \pm 21$ \\
\enddata
\tablecomments{$T_{\rm eq}$ assumes zero Bond albedo and full heat redistribution. $R_p$, $a$, $S$ and $T_{\rm eq}$ use the stellar parameters of Table~\ref{tab:star}.}
\end{deluxetable}

\section{Vetting} \label{sec:vetting}

\subsection{Location of the Signal} \label{sec:centroid}

For each transit we form a difference image (out-of-transit minus in-transit) from the full-frame image cutouts and fit a Gaussian point-spread function to it \citep[cf.][]{Bryson2013, Twicken2018}. Each offset is measured relative to the target's own fitted position in the out-of-transit image, which removes errors in the image's world-coordinate solution and the effect of proper motion. Combining the four transits gives an offset of $1.5'' \pm 1.3''$ from \tic, consistent with the target at $1.2\sigma$. We validated the method on 40 confirmed \tess\ planets: it placed 30 on their host star, left 5 unresolved from a nearby source, returned an offset without a catalogued counterpart for 3, failed for 2, and placed none on a neighbouring star.

Every \gaia\ DR3 source bright enough to produce the observed depth, even if fully eclipsed, lies $21\sigma$ or more from the measured position, with one exception: \gaia\ DR3 5828441289812582656 ($G = 16.15$, $\Delta G = 6.37$), whose parallax of $0.69 \pm 0.08$~mas places it in the background. Owing to the target's proper motion, its separation increases from $3.10''$ at the \gaia\ epoch (2016.0) to $3.4''$--$3.6''$ during the transits; the centroid lies $1.6\sigma$ from its position at the mean transit epoch, so the difference images alone cannot exclude it.

\subsection{Blended Eclipsing Binary} \label{sec:blend}

Diluted by \tic, this star could reproduce the observed depth only through an eclipse of about 74\% of its own light, which requires a companion of comparable size and produces long ingress and egress. We fitted the four transits with eclipse models of the background star for companion-to-star radius ratios $k = 0.3$ to $1.2$, with the impact parameter and $a/R_\star$ free, the dilution free but bounded at 0.4\% of the total flux (40\% above the star's share, to allow for its unknown \tess-band colour), and three limb-darkening laws. Every such model is worse than the planet model by $\Delta\chi^2 \geq 142$ for 2227 data points (Figure~\ref{fig:blend}). A blended eclipsing binary on this star is therefore strongly disfavoured, although high-resolution imaging is required to exclude unresolved sources closer than about $1''$.

\begin{figure}
\plotone{figures/fig4.pdf}
\caption{Diluted-eclipse models for the background star \gaia\ DR3 5828441289812582656, compared with the planet model. Each point is the minimum $\chi^2$ difference over impact parameter, $a/R_\star$ and dilution for a fixed companion radius ratio $k$; symbols denote the assumed limb darkening of the background star.\label{fig:blend}}
\end{figure}

\subsection{Secondary Eclipse and Odd--Even Depths} \label{sec:secondary}

The light curve at orbital phase 0.5, sampled in Sectors~39, 101 and~102, shows no eclipse: the weighted depth is $-54 \pm 39$~ppm, below 62~ppm at $3\sigma$, which excludes an eclipsing companion of comparable surface brightness on a circular orbit. A comparison of odd and even transits is not informative, because only one of the four transits (Sector~100) falls on an odd epoch; the agreement of all four depths (\S\ref{sec:times}) nonetheless disfavours an eclipsing binary at twice the period with unequal eclipses.

\subsection{Statistical False-Positive Probability} \label{sec:fpp}

We computed scenario probabilities with \texttt{TRICERATOPS} \citep{Giacalone2020, Giacalone2021}, using the four transits stacked on their measured mid-times (10-minute bins), a depth of 2100~ppm, the photometric aperture of each sector, and the \gaia\ DR3 field. Ten independent runs give a false-positive probability ${\rm FPP} = 0.003 \pm 0.004$ (largest value 0.012) and a nearby false-positive probability ${\rm NFPP} = 4 \times 10^{-24}$. The planet hypothesis on the target carries 85\% of the probability and a planet diluted by an unresolved bound companion 14\%; the leading false-positive scenario, a planet orbiting an unresolved companion, carries 0.3\%. With the larger default $5 \times 5$-pixel apertures, ${\rm FPP} = 0.001$ and ${\rm NFPP} = 6 \times 10^{-7}$. These values satisfy the numerical criteria of \citet{Giacalone2021}, ${\rm FPP} < 0.015$ and ${\rm NFPP} < 0.001$, but statistical validation also requires high-resolution imaging, which does not yet exist for this star. We therefore classify \planet\ as a planet candidate.

\section{Discussion} \label{sec:discussion}

With $R_p = 3.75 \pm 0.22\,\rearth$, \planet\ lies above the radius valley \citep{Fulton2017}, among sub-Neptunes that likely retain a substantial volatile envelope. At $0.206$~AU it receives ten times Earth's insolation. The mass--radius relation of \citet{ChenKipping2017} predicts $M_p \approx 14\,M_\oplus$ and a radial-velocity semi-amplitude of $K \approx 2.8$~m~s$^{-1}$, an amplitude accessible to current high-precision spectrographs for a $V = 10.0$ star.

Three observations would most efficiently confirm the planet: high-resolution imaging, to exclude unresolved companions and complete the statistical validation; additional transits, to refine the ephemeris and test the timing residuals; and radial velocities, to measure the mass. The transit-timing scatter (\S\ref{sec:ttv}) is itself a target for follow-up, since a nearby perturber would make the system a candidate for dynamical mass measurement.

\section{Predicted Transits} \label{sec:predictions}

Table~\ref{tab:predictions} lists the next transits predicted by the ephemeris of Table~\ref{tab:params}. The quoted uncertainty, $\sigma_{\rm pred} = (\sigma_{T_0}^2 + n^2\sigma_P^2 + \sigma_{\rm rms}^2)^{1/2}$, adds the observed residual scatter $\sigma_{\rm rms} = 21.5$~min to the ephemeris uncertainty, so that it describes an individual future transit rather than the mean ephemeris. \tess\ is expected to observe \tic\ in Sectors~113 and~114 \citep{Burke2020}; the transit of 2027 February~19 is predicted to fall in Sector~114. The transit of 2027 May~3 is observable in full from the southern hemisphere (e.g., La Silla Observatory). A transit absent from the predicted window would falsify the reported period.

\begin{deluxetable}{ccc}
\tablecaption{Predicted transits of \planet \label{tab:predictions}}
\tablehead{\colhead{Epoch} & \colhead{$T_c$ (BJD$_{\rm TDB}$)} & \colhead{$\sigma_{\rm pred}$ (min)}}
\startdata
57 & 2461455.980 & 31 \\
58 & 2461492.320 & 31 \\
59 & 2461528.660 & 31 \\
60 & 2461565.000 & 32 \\
61 & 2461601.340 & 32 \\
62 & 2461637.680 & 32 \\
63 & 2461674.020 & 33 \\
64 & 2461710.359 & 33 \\
65 & 2461746.699 & 34 \\
66 & 2461783.039 & 34 \\
\enddata
\tablecomments{Epochs count orbits from the Sector~39 transit. Epoch~57 (2027 February~19) is expected in \tess\ Sector~114; epoch~59 (2027 May~3, 03:43~UTC at La Silla) is fully observable at night from Chile.}
\end{deluxetable}


\section{Summary} \label{sec:summary}

We report \planet, a $3.75\,\rearth$ sub-Neptune candidate with a 36.34-day period around the nearby K dwarf \tic, detected in four \tess\ full-frame-image transits spanning five years. The period is unique, the transit shape is consistent with the host star's density, the signal is located on the target or within $3.5''$ of it, a blend with the only catalogued neighbour able to mimic it is strongly disfavoured, and the statistical false-positive probability is low. High-resolution imaging and radial velocities are needed to establish its planetary nature, and the transits predicted for 2027 provide a direct test of the ephemeris.

\begin{acknowledgments}
\nolinenumbers
This paper includes data collected by the \tess\ mission, which are publicly available from the Mikulski Archive for Space Telescopes (MAST). Funding for the \tess\ mission is provided by NASA's Science Mission Directorate. This work has made use of data from the European Space Agency (ESA) mission \gaia\ (\url{https://www.cosmos.esa.int/gaia}), processed by the \gaia\ Data Processing and Analysis Consortium (DPAC, \url{https://www.cosmos.esa.int/web/gaia/dpac/consortium}). Funding for the DPAC has been provided by national institutions, in particular the institutions participating in the \gaia\ Multilateral Agreement. This research has made use of the Exoplanet Follow-up Observation Program (ExoFOP; DOI: 10.26134/ExoFOP5) website, which is operated by the California Institute of Technology, under contract with the National Aeronautics and Space Administration under the Exoplanet Exploration Program, and of the NASA Exoplanet Archive, which is operated by the California Institute of Technology under the same program. Software development, data analysis and the drafting of this manuscript were assisted by an AI system (Claude, Anthropic); the author directed the work, verified the results and takes full responsibility for the content.
\end{acknowledgments}

\facilities{TESS, Gaia}

\software{astropy \citep{Astropy2022}, astrocut \citep{Brasseur2019}, batman \citep{Kreidberg2015}, emcee \citep{ForemanMackey2013}, lightkurve \citep{Lightkurve2018}, matplotlib \citep{Hunter2007}, numpy \citep{Harris2020}, scipy \citep{Virtanen2020}, tess-point \citep{Burke2020}, TRICERATOPS \citep{Giacalone2020}}

\section*{Data Availability}
The light curves of all transit windows, our photometry of Sectors~100--103, the measured transit times, the posterior summaries, and the analysis code are available on Zenodo \citep{DataRecord}.

\bibliography{refs}{}
\bibliographystyle{aasjournal}

\end{document}
